\documentclass[11pt]{article}

\usepackage{amsmath, amssymb, amsthm}
\usepackage{algorithm}
\usepackage{algpseudocode}
\usepackage{booktabs}
\usepackage{hyperref}
\usepackage{geometry}
\usepackage{xcolor}
\usepackage{listings}
\usepackage{graphicx}

\geometry{margin=1in}

\newtheorem{theorem}{Theorem}
\newtheorem{lemma}[theorem]{Lemma}
\newtheorem{corollary}[theorem]{Corollary}
\newtheorem{proposition}[theorem]{Proposition}
\newtheorem{definition}{Definition}
\newtheorem{example}{Example}
\newtheorem{remark}{Remark}

\lstset{
  language=C++,
  basicstyle=\ttfamily\small,
  keywordstyle=\color{blue},
  commentstyle=\color{gray},
  breaklines=true,
  frame=single
}

\title{\textbf{A Chromatic Number Algorithm via\\
Inclusion-Exclusion over Independent Sets}}

\author{Graph Exact Coloring Heuristic Project}
\date{\today}

\begin{document}

\maketitle

\begin{abstract}
We describe and analyze an algorithm that determines the chromatic number of a
graph by evaluating a combinatorial quantity derived from inclusion-exclusion
over the collection of all independent sets.  The quantity, denoted $c_k(G)$,
counts the number of ordered $k$-tuples of independent sets of $G$ whose union
is $V(G)$.  We prove that $c_k(G) = 0$ if and only if $k < \chi(G)$, so the
chromatic number equals the smallest $k$ for which $c_k(G) \neq 0$.  We also
describe a companion IS-removal coloring algorithm that computes an upper bound on
$\chi(G)$.  Both algorithms are implemented in C++17 using the Boost Graph
Library.  We present experimental results on cycles $C_n$, complete graphs $K_n$,
the Petersen graph $GP(5,2)$ and other generalized Petersen graphs, wheel
graphs $W_n$, complete multipartite graphs $K_{s,\ldots,s}$, edgeless graphs,
and graphs derived from matrices in the SuiteSparse Matrix Collection
(DNA electrophoresis, graph drawing, and FEM problems), confirming that the
exact algorithm recovers the known chromatic numbers for all tested cases.
We also report empirical scaling data showing that the algorithm is practical
for graphs up to roughly 13--14 vertices on a single core.
\end{abstract}

\tableofcontents
\vspace{1em}

% ─────────────────────────────────────────────────────────────────────────────
\section{Introduction}
% ─────────────────────────────────────────────────────────────────────────────

Graph coloring is one of the most extensively studied problems in combinatorics
and theoretical computer science.  Given a graph $G=(V,E)$, a \emph{proper
$k$-coloring} assigns one of $k$ colors to each vertex so that no two adjacent
vertices share a color.  The minimum such $k$ is the \emph{chromatic number}
$\chi(G)$.  Determining $\chi(G)$ is NP-hard in general~\cite{Karp1972}, yet
exact algorithms are essential for small or sparse graphs arising in scientific
computing, scheduling, and register allocation.

This paper documents the algorithm implemented in the
\texttt{graph-exact-coloring-heuristic} project.  The core function
\texttt{c\_k} evaluates the inclusion-exclusion sum
\[
  c_k(G,k) \;=\; \sum_{S \subseteq V} (-1)^{|S|}\,
                  \bigl|\{\,I \in \mathcal{F}(G) : I \cap S = \emptyset\,\}\bigr|^k,
\]
where $\mathcal{F}(G)$ is the family of all independent sets of $G$ (including
the empty set).  We show that this quantity is zero precisely when $G$ is not
$k$-colorable, and hence one can binary-search (or linearly scan) over $k$ to
find $\chi(G)$.

The companion IS-removal coloring algorithm, implemented in
\texttt{demo\_is\_removal\_coloring.cpp}, repeatedly extracts an independent set
and assigns it the next available color, yielding an upper bound that may
exceed $\chi(G)$ in general.

\paragraph{Organization.}
Section~\ref{sec:prelim} fixes notation and recalls standard facts about graph
coloring.  Section~\ref{sec:algorithm} derives the formula and proves its
correctness.  Section~\ref{sec:heuristic} describes the IS-removal coloring algorithm.
Section~\ref{sec:implementation} outlines the C++ implementation.
Section~\ref{sec:experiments} presents experimental results, including
graphs from the SuiteSparse Matrix Collection.
Section~\ref{sec:complexity} analyzes time and space complexity.
Section~\ref{sec:conclusion} concludes.

% ─────────────────────────────────────────────────────────────────────────────
\section{Preliminaries}\label{sec:prelim}
% ─────────────────────────────────────────────────────────────────────────────

Throughout, $G = (V, E)$ is a simple undirected graph with $n = |V|$ vertices
and $m = |E|$ edges.  Vertices are labeled $0, 1, \ldots, n-1$.

\begin{definition}[Independent set]
  A set $I \subseteq V$ is \emph{independent} if no two vertices in $I$ are
  adjacent.  We write $\mathcal{F}(G)$ for the family of all independent sets
  of $G$, including $\emptyset$.
\end{definition}

\begin{definition}[Proper coloring and chromatic number]
  A \emph{proper $k$-coloring} of $G$ is a function $f: V \to \{1,\ldots,k\}$
  with $f(u) \neq f(v)$ whenever $uv \in E$.  The \emph{chromatic number}
  $\chi(G)$ is the minimum $k$ for which a proper $k$-coloring exists.
\end{definition}

\begin{definition}[Independent set cover]
  An \emph{ordered $k$-cover by independent sets} (hereafter, $k$-IS-cover) is
  a $k$-tuple $(I_1, \ldots, I_k) \in \mathcal{F}(G)^k$ such that
  $I_1 \cup \cdots \cup I_k = V$.
\end{definition}

Note that in a $k$-IS-cover the sets need not be disjoint, and they may be
empty.  A proper $k$-coloring induces a $k$-IS-cover (the color classes), but
a $k$-IS-cover may have overlapping sets.  We will show these notions are
equivalent for the \emph{existence} question.

\begin{definition}[Chromatic polynomial]
  The \emph{chromatic polynomial} $P(G, k)$ counts the number of proper
  $k$-colorings of $G$.  It satisfies $P(G, k) = 0$ iff $k < \chi(G)$.
\end{definition}

Well-known chromatic polynomials for families used in experiments:
\begin{align*}
  P(K_n, k)   &= k(k-1)(k-2)\cdots(k-n+1), \\
  P(C_n, k)   &= (k-1)^n + (-1)^n(k-1),    \\
  P(E_n, k)   &= k^n,
\end{align*}
where $K_n$ is the complete graph, $C_n$ the cycle, and $E_n$ the edgeless
graph on $n$ vertices.

% ─────────────────────────────────────────────────────────────────────────────
\section{The Exact Algorithm}\label{sec:algorithm}
% ─────────────────────────────────────────────────────────────────────────────

\subsection{The Combinatorial Formula}

For a graph $G$ and a positive integer $k$, define
\[
  \alpha(S,\,\mathcal{F}) \;=\;
    \bigl|\{\,I \in \mathcal{F}(G) : I \cap S = \emptyset\,\}\bigr|
\]
to be the number of independent sets of $G$ that avoid $S$ entirely.
Equivalently, $\alpha(S,\mathcal{F})$ is the total number of independent sets
in the subgraph of $G$ induced by $V \setminus S$.  The algorithm evaluates
\begin{equation}\label{eq:ck}
  c_k(G, k)
  \;=\;
  \sum_{S \subseteq V} (-1)^{|S|}\,\alpha(S,\,\mathcal{F})^k.
\end{equation}

\subsection{Combinatorial Interpretation}

\begin{theorem}\label{thm:ck-cover}
  $c_k(G,k)$ equals the number of ordered $k$-IS-covers of $G$.
\end{theorem}

\begin{proof}
Let $N = |\mathcal{F}(G)|$.  The number of ordered $k$-tuples
$(I_1,\ldots,I_k) \in \mathcal{F}(G)^k$ with no constraint is $N^k$.  For a
fixed vertex $v \in V$ let $A_v$ be the set of $k$-tuples in which $v$ is
covered by no $I_j$ (i.e.\ $v \notin I_j$ for all $j$).  Then
\[
  |A_v| = |\{I \in \mathcal{F}(G) : v \notin I\}|^k = \alpha(\{v\},\mathcal{F})^k.
\]
More generally, for any $S \subseteq V$,
\[
  \Bigl|\bigcap_{v \in S} A_v\Bigr|
  = |\{I \in \mathcal{F}(G) : S \cap I = \emptyset\}|^k
  = \alpha(S, \mathcal{F})^k.
\]
By inclusion-exclusion, the number of $k$-tuples in which every vertex is
covered by at least one $I_j$ is
\[
  \sum_{S \subseteq V} (-1)^{|S|}\,\alpha(S,\mathcal{F})^k
  \;=\; c_k(G,k). \qedhere
\]
\end{proof}

\subsection{Equivalence with $k$-Colorability}

\begin{theorem}\label{thm:chromatic}
  $c_k(G, k) = 0$ if and only if $G$ is not properly $k$-colorable,
  i.e.\ $k < \chi(G)$.
\end{theorem}

\begin{proof}
$(\Leftarrow)$ Suppose $k \geq \chi(G)$.  Fix a proper $k$-coloring
$f: V \to \{1,\ldots,k\}$ and let $I_j = f^{-1}(j)$ for each $j$.  Each $I_j$
is an independent set and $\bigcup_j I_j = V$, so $(I_1,\ldots,I_k)$ is a
$k$-IS-cover.  Hence $c_k(G,k) \geq 1 > 0$.

$(\Rightarrow)$ Suppose $k < \chi(G)$.  We show no $k$-IS-cover exists.
If $(I_1,\ldots,I_k)$ were a $k$-IS-cover, define $f(v) = \min\{j : v \in I_j\}$
for each $v \in V$.  Since $I_j$ is independent, no two adjacent vertices can
share the same $j$; thus $f$ is a proper $k$-coloring, contradicting
$k < \chi(G)$.  Therefore $c_k(G,k) = 0$.
\end{proof}

\begin{corollary}
  $\chi(G) = \min\{\,k \geq 1 : c_k(G,k) \neq 0\,\}.$
\end{corollary}

\begin{remark}
  The quantity $c_k(G,k)$ is not the chromatic polynomial $P(G,k)$ in general.
  For $k = \chi(G)$, the two coincide for complete graphs (where every
  $k$-IS-cover induces a proper coloring immediately), but they differ for
  most other graphs at values $k > \chi(G)$.  For example, $c_k(C_5, 3) = 60$
  while $P(C_5, 3) = 30$.  Nevertheless, both are zero for $k < \chi(G)$ and
  nonzero for $k \geq \chi(G)$, which is the only property needed to determine
  the chromatic number.
\end{remark}

\subsection{Algorithm}

The full procedure to compute $\chi(G)$ is given in
Algorithm~\ref{alg:exact}.  It first enumerates all independent sets of $G$
(Algorithm~\ref{alg:indsets}) using a power-set sweep with bitset membership
testing, then evaluates formula~\eqref{eq:ck} for increasing $k$.

\begin{algorithm}
\caption{Enumerate all independent sets of $G$}\label{alg:indsets}
\begin{algorithmic}[1]
\Require Graph $G = (V, E)$ with $|V| = n$
\Ensure Set $\mathcal{F}$ of all independent sets (as $n$-bit integers)
\State $\mathcal{F} \gets \emptyset$
\For{$i \gets 0$ \textbf{to} $2^n - 1$}
  \State $B \gets$ $n$-bit representation of $i$
  \If{$\forall\, u \neq v$: $B[u] = B[v] = 1 \Rightarrow uv \notin E$}
    \State $\mathcal{F} \gets \mathcal{F} \cup \{B\}$
  \EndIf
\EndFor
\Return $\mathcal{F}$
\end{algorithmic}
\end{algorithm}

\begin{algorithm}
\caption{Compute $\chi(G)$ via inclusion-exclusion}\label{alg:exact}
\begin{algorithmic}[1]
\Require Graph $G = (V, E)$ with $|V| = n$
\Ensure Chromatic number $\chi(G)$
\State $\mathcal{F} \gets \textsc{EnumerateIndepSets}(G)$
\For{$k \gets 1$ \textbf{to} $n$}
  \State $\mathit{cnt} \gets 0$
  \For{each $S \subseteq V$ (represented as a bitset)}
    \State $a \gets |\{I \in \mathcal{F} : I \mathbin{\&} S = 0\}|$
      \Comment{count IS avoiding $S$}
    \State $\mathit{cnt} \gets \mathit{cnt} + (-1)^{|S|} \cdot a^k$
  \EndFor
  \If{$\mathit{cnt} \neq 0$}
    \Return $k$
  \EndIf
\EndFor
\end{algorithmic}
\end{algorithm}

\subsubsection*{Bitset Representation}

Each independent set is stored as an $n$-bit integer (a
\texttt{boost::dynamic\_bitset}).  The independence check for a candidate set
$B$ iterates over all pairs $(i,j)$ with $B[i]=B[j]=1$ and verifies $ij \notin
E$.  The subset-avoidance test $I \cap S = \emptyset$ reduces to the bitwise
operation $(I \mathbin{\&} S) = 0$.

% ─────────────────────────────────────────────────────────────────────────────
\section{IS-Removal Coloring}\label{sec:heuristic}
% ─────────────────────────────────────────────────────────────────────────────

\subsection{Independent Set Extraction}

The IS extraction procedure (Algorithm~\ref{alg:greedyIS}) builds one
independent set by iterating through the active vertices in index order.
For each vertex $v$ selected, it removes $v$ and all of $v$'s current
neighbors from the active set.

\begin{algorithm}
\caption{Independent set extraction}\label{alg:greedyIS}
\begin{algorithmic}[1]
\Require Graph $G = (V, E)$; active set $U \subseteq V$
\Ensure An independent set $I \subseteq U$
\State $I \gets \emptyset$
\While{$G$ has edges}
  \State $v \gets$ first vertex in $U$
  \State $I \gets I \cup \{v\}$
  \State Remove $v$ and all neighbors of $v$ from $U$ and from $G$
\EndWhile
\Return $I$
\end{algorithmic}
\end{algorithm}

\subsection{IS-Removal Coloring}

Algorithm~\ref{alg:greedycol} repeatedly calls the IS extractor to peel off
one independent set (color class) at a time.

\begin{algorithm}
\caption{IS-removal coloring}\label{alg:greedycol}
\begin{algorithmic}[1]
\Require Graph $G = (V, E)$
\Ensure Coloring array $\mathit{col}[0\ldots n-1]$
\State $\mathit{col} \gets$ array of size $n$; $k \gets 0$
\While{$G$ has edges}
  \State $I \gets \textsc{ExtractIndepSet}(G)$
  \For{each $v \in I$}
    \State $\mathit{col}[v] \gets k$
    \State Remove all edges incident to $v$ from $G$
  \EndFor
  \State $k \gets k + 1$
\EndWhile
\Return $\mathit{col}$
\end{algorithmic}
\end{algorithm}

\subsection{Quality and Complexity}

IS-removal coloring produces at most $\Delta(G)+1$ colors, where $\Delta(G)$
is the maximum degree~\cite{Welsh1967}.  The vertex-order traversal makes it
exact on complete graphs and odd cycles but can be arbitrarily far from optimal
on bipartite or dense multipartite graphs (see Table~\ref{tab:greedy} for
empirical comparisons).  Its primary advantage is $O(n + m)$ running time per
color class, making it practical for large sparse graphs.

% ─────────────────────────────────────────────────────────────────────────────
\section{Implementation}\label{sec:implementation}
% ─────────────────────────────────────────────────────────────────────────────

The project is written in C++17 and uses the Boost Graph Library (BGL)
version~1.60 or later.  The graph type is

\begin{lstlisting}
typedef boost::adjacency_list<
    boost::vecS, boost::vecS, boost::undirectedS,
    property<boost::vertex_color_t, int>,
    property<boost::edge_weight_t, int,
             property<boost::edge_name_t, string>>> Graph;
\end{lstlisting}

Independent sets are stored as \texttt{boost::dynamic\_bitset<>} objects, which
support $O(n/w)$ bitwise AND, OR, and population-count operations where $w$ is
the machine word size.

The key functions are:

\begin{itemize}
  \item \texttt{power\_set(n, f)}: iterates over all $2^n$ subsets, calling
        \texttt{f} with each subset represented as a \texttt{dynbit}.
  \item \texttt{is\_indep\_set(G, B)}: checks all pairs of set bits in $B$ for
        adjacency.
  \item \texttt{gen\_ind\_set\_for\_g(G)}: returns $\mathcal{F}(G)$ as a
        \texttt{std::set<dynbit>}.
  \item \texttt{alpha(S, F)}: counts $|\{I \in F : I \mathbin{\&} S = 0\}|$
        by a linear scan.
  \item \texttt{c\_k(F, G, k)}: evaluates formula~\eqref{eq:ck} using
        floating-point arithmetic and casts the result to \texttt{int}.
\end{itemize}

Generators for Cycle ($C_n$) and Generalized Peterson ($GP(n,k)$) graphs
extend the abstract base class \texttt{GeneratorInterface} and override
\texttt{generate()}.

\paragraph{Generalized Peterson graph.}
$GP(n,k)$ has $2n$ vertices.  The outer $n$-cycle connects vertex $i$ to
$(i+1) \bmod n$.  Spoke edges connect outer vertex $i$ to inner vertex $i+n$.
Inner edges connect $i+n$ to $(i+n+k) \bmod n + n$ with a boundary correction.

\paragraph{Build system.}
The project uses CMake 3.6 or later.  The \texttt{exact\_coloring} target
compiles \texttt{main.cpp} (exact algorithm demo) and the \texttt{tests} target
compiles \texttt{tests.cpp} (the test suite described in
Section~\ref{sec:experiments}).

% ─────────────────────────────────────────────────────────────────────────────
\section{Experimental Results}\label{sec:experiments}
% ─────────────────────────────────────────────────────────────────────────────

All experiments were run on the \texttt{tests} binary compiled with
\texttt{-std=c++17 --openmp -O2} on Ubuntu 22.04.  Table~\ref{tab:results}
summarizes the chromatic numbers recovered by the exact algorithm.

\begin{table}[h]
\centering
\caption{Chromatic numbers computed by the exact algorithm (78 assertions, all passing).}
\label{tab:results}
\begin{tabular}{llcccc}
\toprule
Graph & Description & $|V|$ & $|E|$ & $\chi$ (computed) & Known $\chi$ \\
\midrule
\multicolumn{6}{l}{\textit{Cycle graphs}} \\
$C_3$ & Triangle / $K_3$         & 3  & 3  & 3 & 3 \\
$C_4$ & 4-cycle (square)         & 4  & 4  & 2 & 2 \\
$C_5$ & 5-cycle (pentagon)       & 5  & 5  & 3 & 3 \\
$C_6$ & 6-cycle (hexagon)        & 6  & 6  & 2 & 2 \\
$C_7$ & 7-cycle                  & 7  & 7  & 3 & 3 \\
\midrule
\multicolumn{6}{l}{\textit{Complete and edgeless graphs}} \\
$K_1$ & Single vertex            & 1  & 0  & 1 & 1 \\
$K_2$ & Single edge              & 2  & 1  & 2 & 2 \\
$K_3$ & Triangle                 & 3  & 3  & 3 & 3 \\
$K_4$ & Complete graph           & 4  & 6  & 4 & 4 \\
$K_9$ & Complete graph           & 9  & 36 & 9 & 9 \\
$K_{10}$ & Complete graph        & 10 & 45 & 10 & 10 \\
$E_1$ & Isolated vertex          & 1  & 0  & 1 & 1 \\
$E_3$ & Three isolated vertices  & 3  & 0  & 1 & 1 \\
\midrule
\multicolumn{6}{l}{\textit{Bipartite graphs}} \\
$K_{2,3}$ & Complete bipartite   & 5  & 6  & 2 & 2 \\
\midrule
\multicolumn{6}{l}{\textit{Wheel graphs} $W_n$ = hub $+\,C_{n-1}$;
  $\chi=3$ if $n$ odd, $\chi=4$ if $n$ even} \\
$W_5$ & Hub $+\,C_4$ (even rim)     & 5  & 8  & 3 & 3 \\
$W_6$ & Hub $+\,C_5$ (odd rim)      & 6  & 10 & 4 & 4 \\
$W_7$ & Hub $+\,C_6$ (even rim)     & 7  & 12 & 3 & 3 \\
$W_8$ & Hub $+\,C_7$ (odd rim)      & 8  & 14 & 4 & 4 \\
$W_9$ & Hub $+\,C_8$ (even rim)     & 9  & 16 & 3 & 3 \\
\midrule
\multicolumn{6}{l}{\textit{Complete multipartite graphs};
  $\chi(K_{s,\ldots,s}^{(k)}) = k$} \\
$K_{2,2,2}$   & Octahedron            & 6  & 12 & 3 & 3 \\
$K_{3,3,3}$   & Complete 3-partite    & 9  & 27 & 3 & 3 \\
$K_{2,2,2,2}$ & Complete 4-partite    & 8  & 24 & 4 & 4 \\
$K_{2\times5}$ & Complete 5-partite   & 10 & 40 & 5 & 5 \\
$K_{2\times6}$ & Complete 6-partite   & 12 & 60 & 6 & 6 \\
$K_{2\times7}$ & Complete 7-partite   & 14 & 84 & 7 & 7 \\
$K_{2\times8}$ & Complete 8-partite   & 16 & 112 & 8 & 8 \\
$K_{2\times9}$ & Complete 9-partite   & 18 & 144 & 9 & 9 \\
\midrule
\multicolumn{6}{l}{\textit{Generalized Petersen graphs} $GP(n,k)$} \\
$GP(4,1)$ & 3-cube $Q_3$ (bipartite)    & 8  & 12 & 2 & 2 \\
$GP(5,1)$ & Pentagonal prism            & 10 & 15 & 3 & 3 \\
$GP(5,2)$ & Petersen graph              & 10 & 15 & 3 & 3 \\
$GP(6,1)$ & Hexagonal prism (bipartite) & 12 & 18 & 2 & 2 \\
$GP(6,2)$ & ---                         & 12 & 18 & 3 & 3 \\
$GP(7,2)$ & ---                         & 14 & 21 & 3 & 3 \\
\bottomrule
\end{tabular}
\end{table}

All 78 test assertions pass, confirming both the zero/nonzero behavior of
$c_k$ and, for complete graphs, the exact value of the formula at $k = \chi(G)$.
In particular, $c_k(K_9, 9) = 9! = 362{,}880$ and $c_k(K_{10}, 10) = 10! = 3{,}628{,}800$
are verified exactly, confirming the correctness of the \texttt{long long} accumulation.

\subsection{Exact Values of $c_k$ at the Chromatic Number}

For complete graphs, the formula~\eqref{eq:ck} coincides with the standard
chromatic polynomial at $k = \chi(G)$:

\begin{table}[h]
\centering
\caption{Values of $c_k(G, \chi(G))$ vs.\ the chromatic polynomial $P(G, \chi(G))$.}
\label{tab:poly}
\begin{tabular}{lccccc}
\toprule
Graph & $\chi$ & $c_k(G, \chi)$ & $P(G, \chi)$ & Match? \\
\midrule
$K_1$ & 1 & 1    & 1    & Yes \\
$K_2$ & 2 & 2    & 2    & Yes \\
$K_3$ & 3 & 6    & 6    & Yes \\
$K_4$ & 4 & 24   & 24   & Yes \\
$C_3$ & 3 & 6    & 6    & Yes \\
$C_4$ & 2 & 2    & 2    & Yes \\
$C_5$ & 3 & 60   & 30   & No  \\
\bottomrule
\end{tabular}
\end{table}

The coincidence for $K_n$ can be verified algebraically: since the only
independent sets of $K_n$ are $\emptyset$ and the $n$ singletons,
$\alpha(S, \mathcal{F}) = n - |S| + 1$ and the sum~\eqref{eq:ck} at $k=n$
evaluates to
\[
  \sum_{j=0}^{n} \binom{n}{j}(-1)^{n-j}(j+1)^n = n!,
\]
which equals $P(K_n, n) = n!$.  For cycle $C_4$ the coincidence is accidental;
for $C_5$ the values differ.

\subsection{Parity Property of Odd Cycles}

For any odd cycle $C_n$ ($n$ odd) the algorithm correctly returns $\chi = 3$.
No odd cycle is 2-colorable because any 2-coloring of a path forces the two
endpoints of a cycle to receive the same color, contradicting the closing edge.
For $n = 3, 5, 7$ the algorithm confirmed $c_k(C_n, 2) = 0$ and $c_k(C_n, 3)
> 0$.

\subsection{Petersen Graph}

The Petersen graph $GP(5,2)$ is the canonical example of a 3-chromatic
3-regular graph.  It is not planar and has no Hamiltonian cycle.  Our algorithm
confirmed $\chi(GP(5,2)) = 3$, with $c_k(GP(5,2), 2) = 0$ and $c_k(GP(5,2),
3) > 0$.

\subsection{Wheel Graphs}

The wheel graph $W_n$ consists of a single hub vertex connected to every vertex
of a cycle $C_{n-1}$.  It has $n$ vertices and $2(n-1)$ edges.  Because the
hub is adjacent to every rim vertex, the hub's color must differ from all colors
used on the rim.

\begin{itemize}
  \item If $n$ is odd, the rim $C_{n-1}$ has an even number of vertices and is
        therefore bipartite, so $\chi(C_{n-1}) = 2$.  The hub requires a third
        color, giving $\chi(W_n) = 3$.
  \item If $n$ is even, the rim $C_{n-1}$ is an odd cycle, so $\chi(C_{n-1})
        = 3$.  The hub must avoid all three rim colors, giving $\chi(W_n) = 4$.
\end{itemize}

This makes wheel graphs the simplest natural family with $\chi = 4$.  We
verified $\chi(W_n)$ for $n \in \{5,6,7,8,9\}$; all 5 tests pass.

\subsection{Complete Multipartite Graphs}

The complete $k$-partite graph $K_{s,\ldots,s}$ (with $k$ parts of $s$
vertices each) has $\chi = k$: color each part with its own color (lower bound
from the clique-like independence structure; $k$ colors suffice since each part
is an independent set).  We tested:

\begin{itemize}
  \item $K_{2,2,2}$ (octahedron): 6 vertices, 12 edges, $\chi = 3$.
  \item $K_{3,3,3}$: 9 vertices, 27 edges, $\chi = 3$.
  \item $K_{2,2,2,2}$: 8 vertices, 24 edges, $\chi = 4$.
\end{itemize}

These graphs have few independent sets relative to $2^n$ (only subsets within
a single part), so the algorithm is fast despite the large edge count.

\subsection{Generalized Petersen Family}

Beyond $GP(5,2)$ we tested five more members:

\begin{itemize}
  \item $GP(4,1) \cong Q_3$ (3-cube): bipartite, $\chi = 2$.
  \item $GP(5,1)$ (pentagonal prism): odd outer cycle, $\chi = 3$.
  \item $GP(6,1)$ (hexagonal prism): even outer cycle, bipartite, $\chi = 2$.
  \item $GP(6,2)$: inner triangles force $\chi \geq 3$; algorithm confirms $\chi = 3$.
  \item $GP(7,2)$: odd outer and inner cycles, $\chi = 3$.
\end{itemize}

\subsection{SuiteSparse Matrix Collection Graphs}\label{sec:suitesparse}

To demonstrate applicability to real-world data, we obtained matrices from
the SuiteSparse Matrix Collection~\cite{SuiteSparse} and converted each to
an undirected graph by symmetrizing (adding edge $(i,j)$ whenever $A_{ij}
\neq 0$ or $A_{ji} \neq 0$) and discarding self-loops and edge weights.
For matrices larger than 14 vertices we extract a connected subgraph of $k$
vertices via BFS from a fixed starting vertex.

\begin{table}[h]
\centering
\caption{Chromatic numbers on graphs derived from SuiteSparse matrices.}
\label{tab:suitesparse}
\begin{tabular}{llllcccc}
\toprule
Matrix & Group & Domain & Processing & $|V|$ & $|E|$ & $\chi$ & Time \\
\midrule
\texttt{cage3}   & vanHeukelum & DNA electrophoresis & symmetrized (full) & 5  &  7 & 3 & $<$1\,ms \\
\texttt{cage4}   & vanHeukelum & DNA electrophoresis & symmetrized (full) & 9  & 20 & 3 &    3\,ms \\
\texttt{GD98\_a} & Pajek       & Graph Drawing 1998  & BFS-10 subgraph    & 10 & 10 & 3 &   81\,ms \\
\texttt{ndtest}  & (project)   & FEM/structural      & BFS-10 subgraph    & 10 & 14 & 2 &   15\,ms \\
\bottomrule
\end{tabular}
\end{table}

The cage matrices arise from random walks on finite words over the DNA
alphabet $\{A, C, G, T\}$; \texttt{cage3} has $|\Sigma|^3 - |\Sigma|^2 = 48$
directed edges reduced to 7 undirected after symmetrization.  Both DNA graphs
require exactly 3 colors --- they contain odd cycles formed by overlapping
DNA words that share a common suffix/prefix.  The \texttt{GD98\_a} subgraph
also requires 3 colors, while the symmetric FEM matrix \texttt{ndtest} is
2-colorable in its 10-vertex neighborhood, consistent with the bipartite
structure typical of finite-element adjacency graphs.

\subsection{Scaling Analysis}

Table~\ref{tab:scaling} reports wall-clock times for selected larger graphs.
The dominant cost is evaluating $c_k$ (formula~\eqref{eq:ck}), which is
$O(2^n \cdot |\mathcal{F}|)$ per candidate $k$.  For sparse graphs such as
cycles, $|\mathcal{F}|$ grows like the Lucas sequence $L(n) \approx \varphi^n$,
so the total cost is roughly $O((2\varphi)^n) \approx O(3.24^n)$, yielding
approximately a 10$\times$ slowdown per 2 additional vertices.  Dense graphs
such as $K_{m,m}$ have much smaller $|\mathcal{F}|$ and therefore scale better.

\begin{table}[h]
\centering
\caption{Wall-clock time for larger graphs (Release build, single core, Ubuntu 22.04).}
\label{tab:scaling}
\begin{tabular}{llcccc}
\toprule
Graph & Family & $|V|$ & $|E|$ & $\chi$ & Time \\
\midrule
$C_8$     & Cycle          &  8 &  8 & 2 &    1\,ms \\
$C_9$     & Cycle          &  9 &  9 & 3 &    7\,ms \\
$C_{10}$  & Cycle          & 10 & 10 & 2 &   24\,ms \\
$C_{11}$  & Cycle          & 11 & 11 & 3 &   86\,ms \\
$C_{12}$  & Cycle          & 12 & 12 & 2 &  168\,ms \\
$C_{13}$  & Cycle          & 13 & 13 & 3 &  786\,ms \\
$C_{14}$  & Cycle          & 14 & 14 & 2 & $\approx$3\,s \\
$C_{16}$  & Cycle          & 16 & 16 & 2 & $\approx$33\,s \\
\midrule
$K_5$     & Complete       &  5 & 10 & 5 & $<$1\,ms \\
$K_6$     & Complete       &  6 & 15 & 6 & $<$1\,ms \\
$K_7$     & Complete       &  7 & 21 & 7 & $<$1\,ms \\
$K_8$     & Complete       &  8 & 28 & 8 &    1\,ms \\
\midrule
$K_{4,4}$ & Bipartite      &  8 & 16 & 2 &    1\,ms \\
$K_{5,5}$ & Bipartite      & 10 & 25 & 2 &    9\,ms \\
$K_{6,6}$ & Bipartite      & 12 & 36 & 2 &   69\,ms \\
$K_{7,7}$ & Bipartite      & 14 & 49 & 2 &  536\,ms \\
\midrule
$GP(6,2)$ & Gen.\ Petersen & 12 & 18 & 3 &  135\,ms \\
$GP(7,2)$ & Gen.\ Petersen & 14 & 21 & 3 & 1333\,ms \\
$GP(8,3)$ & Gen.\ Petersen & 16 & 24 & 2 & $\approx$15\,s \\
\bottomrule
\end{tabular}
\end{table}

\subsection{High Chromatic Number Graphs}

Complete graphs $K_n$ have $\chi(K_n) = n$ and $|\mathcal{F}| = n+1$ (only
singletons and the empty set), so the independent-set enumeration phase is
trivial.  The bottleneck is computing the chromatic polynomial at each candidate
$k$, which requires intermediate terms of the form
$(-1)^{|S|}\,\alpha(S,\mathcal{F})^k$.  For $K_9$, these terms reach
$9^9 \approx 3.9 \times 10^8$, which fits comfortably in a 64-bit integer.
The implementation uses \texttt{long long} accumulation with integer
exponentiation (no floating point), giving correct results up to at least
$K_{10}$ ($10! = 3{,}628{,}800$, verified by the test suite).

Complete $k$-partite graphs $K_{2 \times k}$ (two vertices per part) provide
compact benchmarks for high chromatic numbers: only $2k$ vertices but $\chi = k$.
Their independent-set families have exactly $3^k$ members (each part contributes
three choices: both, neither, or one vertex), far fewer than $2^{2k}$, so the
algorithm scales well despite the large chromatic number.

\begin{table}[h]
\centering
\caption{High chromatic number graphs (Release build, single core).}
\label{tab:highchi}
\begin{tabular}{lcccc}
\toprule
Graph & $V$ & $E$ & $\chi$ & Time \\
\midrule
$K_9$         & 9  & 36 & 9 & 3\,ms \\
$K_{10}$      & 10 & 45 & 10 & 9\,ms \\
$K_{2\times5}$ & 10 & 40 & 5 & 7\,ms \\
$K_{2\times6}$ & 12 & 60 & 6 & 47\,ms \\
$K_{2\times7}$ & 14 & 84 & 7 & 287\,ms \\
$K_{2\times8}$ & 16 & 112 & 8 & 1.1\,s \\
$K_{2\times9}$ & 18 & 144 & 9 & 5.5\,s \\
\bottomrule
\end{tabular}
\end{table}

\subsection{IS-Removal vs.\ Exact Coloring}\label{sec:greedy_vs_exact}

Table~\ref{tab:greedy} compares the number of colors produced by IS-removal
coloring (Algorithm~\ref{alg:greedycol}) against the true chromatic number
$\chi(G)$ on the same graph families.  The algorithm processes vertices in
ascending index order, so its quality depends strongly on the vertex ordering.

\begin{table}[ht]
\centering
\caption{IS-removal coloring vs.\ exact chromatic number $\chi(G)$.
  ``='' means IS-removal is optimal; ``$>$'' means it over-estimates.}
\label{tab:greedy}
\begin{tabular}{llcccc}
\toprule
Graph & Family & $\chi$ & IS-removal & $\Delta(G)$ & Note \\
\midrule
\multicolumn{6}{l}{\textit{Cycles}} \\
$C_3$  & Odd cycle  & 3 & 3  & 2 & = optimal \\
$C_4$  & Even cycle & 2 & 4  & 2 & $2\times$ optimum \\
$C_5$  & Odd cycle  & 3 & 3  & 2 & = optimal \\
$C_6$  & Even cycle & 2 & 4  & 2 & $2\times$ optimum \\
$C_7$  & Odd cycle  & 3 & 3  & 2 & = optimal \\
\midrule
\multicolumn{6}{l}{\textit{Complete graphs}} \\
$K_2$  & Complete & 2  & 2  & 1 & = optimal \\
$K_3$  & Complete & 3  & 3  & 2 & = optimal \\
$K_4$  & Complete & 4  & 4  & 3 & = optimal \\
$K_5$  & Complete & 5  & 5  & 4 & = optimal \\
$K_{10}$ & Complete & 10 & 10 & 9 & = optimal \\
\midrule
\multicolumn{6}{l}{\textit{Bipartite}} \\
$K_{2,3}$ & Complete bipartite & 2 & 3 & 3 & over by 1 \\
\midrule
\multicolumn{6}{l}{\textit{Wheel graphs} $W_n$} \\
$W_5$  & Odd $n$, even rim & 3 & 5 & 4 & over by 2 \\
$W_6$  & Even $n$, odd rim & 4 & 4 & 5 & = optimal \\
$W_7$  & Odd $n$, even rim & 3 & 5 & 6 & over by 2 \\
$W_8$  & Even $n$, odd rim & 4 & 4 & 7 & = optimal \\
\midrule
\multicolumn{6}{l}{\textit{Complete multipartite}} \\
$K_{2,2,2}$   & Octahedron       & 3 & 5  & 4 & over by 2 \\
$K_{2,2,2,2}$ & 4-partite        & 4 & 7  & 6 & over by 3 \\
$K_{2\times5}$ & 5-partite       & 5 & 9  & 8 & over by 4 \\
$K_{2\times9}$ & 9-partite       & 9 & 17 & 16 & over by 8 \\
\midrule
\multicolumn{6}{l}{\textit{Generalized Petersen graphs}} \\
$GP(5,1)$ & Pentagonal prism & 3 & 3 & 3 & = optimal \\
$GP(5,2)$ & Petersen graph   & 3 & 4 & 3 & over by 1 \\
$GP(6,1)$ & Hexagonal prism  & 2 & 6 & 3 & $3\times$ optimum \\
$GP(7,2)$ & ---              & 3 & 4 & 3 & over by 1 \\
\bottomrule
\end{tabular}
\end{table}

The results illustrate several patterns:

\begin{itemize}
\item \textbf{Complete graphs}: IS-removal is always optimal.  Because every
  vertex is adjacent to all lower-indexed vertices, each vertex receives a new
  color, achieving exactly $n$ colors for $K_n$.

\item \textbf{Odd cycles}: IS-removal is optimal ($\chi = 3$).  Vertex 0 and
  its non-adjacent neighbor both receive color 0, and the single remaining
  vertex requires a third color.

\item \textbf{Even cycles and bipartite graphs}: IS-removal performs poorly
  because the vertex-order traversal fails to exploit the bipartite structure.
  $C_4$ receives 4 colors instead of 2; similarly $GP(6,1)$ (hexagonal prism,
  bipartite) receives 6 instead of 2.

\item \textbf{Wheel graphs}: exact when $n$ is even (odd rim forces 4 colors
  regardless of ordering) but over-estimates by 2 when $n$ is odd.

\item \textbf{Complete multipartite $K_{2\times k}$}: the over-count grows
  linearly with $k$; IS-removal uses approximately $2k-1$ colors instead of
  $k$, giving an approximation ratio approaching 2 as $k$ grows.
\end{itemize}

These results confirm that IS-removal coloring is a fast but unreliable upper
bound: it is tight on graphs where the natural vertex order aligns with the
chromatic structure (complete graphs, odd cycles), but can be far from optimal
on structured sparse or multipartite graphs.

% ─────────────────────────────────────────────────────────────────────────────
\section{Complexity Analysis}\label{sec:complexity}
% ─────────────────────────────────────────────────────────────────────────────

\subsection{Independent Set Enumeration}

\texttt{gen\_ind\_set\_for\_g} calls \texttt{power\_set}, which iterates over
all $2^n$ subsets.  For each subset $B$, \texttt{is\_indep\_set} checks all
pairs of set bits: at most $O(n^2)$ edge queries.  Hence the enumeration phase
runs in $O(2^n \cdot n^2)$ time and produces a set $\mathcal{F}$ of at most
$3^n / \text{poly}(n)$ independent sets (by the Moon--Moser theorem on the
maximum number of maximal independent sets in a graph with $n$
vertices~\cite{MoonMoser1965}).

\subsection{Evaluation of $c_k$}

For each candidate $k$, evaluating formula~\eqref{eq:ck} requires iterating
over all $2^n$ subsets $S$ and, for each $S$, scanning all $|\mathcal{F}|$
independent sets to compute $\alpha(S, \mathcal{F})$.  This costs
$O(2^n \cdot |\mathcal{F}| \cdot n/w)$ per value of $k$, where $n/w$ is the
cost of a bitwise AND.  Summing over at most $n$ values of $k$, the total
evaluation cost is $O(n \cdot 2^n \cdot |\mathcal{F}| \cdot n/w)$.

\subsection{Overall Complexity}

The dominating term is the evaluation of $c_k$, giving
\[
  T(n) = O\!\left(n^2 \cdot 2^n \cdot |\mathcal{F}(G)|\right).
\]
In the worst case (e.g., the empty graph $E_n$ has $|\mathcal{F}| = 2^n$),
this is $O(n^2 \cdot 4^n)$.  Empirically (Table~\ref{tab:scaling}), sparse
graphs such as cycles reach the practical limit around $n = 13$--$14$ (seconds
to minutes), while dense graphs such as $K_{m,m}$ scale better because
$|\mathcal{F}|$ is much smaller than $2^n$.  The algorithm is intended for
exact coloring of small subgraphs; for large graphs the IS-removal coloring in
\texttt{demo\_is\_removal\_coloring.cpp} is the appropriate choice.

\subsection{Space Complexity}

Storing $\mathcal{F}$ requires $O(|\mathcal{F}| \cdot n/w)$ words, which is
$O(3^n \cdot n / w)$ in the worst case.

% ─────────────────────────────────────────────────────────────────────────────
\section{Conclusion}\label{sec:conclusion}
% ─────────────────────────────────────────────────────────────────────────────

We have described and formally justified the inclusion-exclusion algorithm for
computing the chromatic number of a graph implemented in the
\texttt{graph-exact-coloring-heuristic} project.  The key result is that the
quantity $c_k(G, k)$ --- the number of ordered $k$-tuples of independent sets
covering $V$ --- is zero if and only if $G$ is not $k$-colorable
(Theorem~\ref{thm:chromatic}).  Scanning $k = 1, 2, \ldots$ and stopping at
the first nonzero value thus yields $\chi(G)$ exactly.

The algorithm is exponential in $|V|$ and is therefore intended for small
graphs or for exact coloring of subgraphs within a larger heuristic pipeline.
The companion IS-removal coloring algorithm runs in $O(n + m)$ time per color
class and provides a fast upper bound useful as an initialization strategy or a
standalone approximation for large instances.

Future directions include: (i) exploiting symmetry to prune the power-set
enumeration, (ii) integrating branch-and-bound to improve worst-case behavior,
and (iii) leveraging OpenMP (already supported in the build) to parallelize the
inner $\alpha$-counting loop.

\begin{thebibliography}{9}

\bibitem{Karp1972}
R.~M. Karp.
\newblock Reducibility among combinatorial problems.
\newblock In R.~E. Miller and J.~W. Thatcher, editors, \emph{Complexity of
  Computer Computations}, pages 85--103. Plenum Press, 1972.

\bibitem{Welsh1967}
D.~J.~A. Welsh and M.~B. Powell.
\newblock An upper bound for the chromatic number of a graph and its
  application to timetabling problems.
\newblock \emph{The Computer Journal}, 10(1):85--86, 1967.

\bibitem{MoonMoser1965}
J.~W. Moon and L.~Moser.
\newblock On cliques in graphs.
\newblock \emph{Israel Journal of Mathematics}, 3(1):23--28, 1965.

\bibitem{Lawler1976}
E.~L. Lawler.
\newblock A note on the complexity of the chromatic number problem.
\newblock \emph{Information Processing Letters}, 5(3):66--67, 1976.

\bibitem{BGL}
J.~G. Siek, L.-Q. Lee, and A.~Lumsdaine.
\newblock \emph{The Boost Graph Library: User Guide and Reference Manual}.
\newblock Addison-Wesley, 2002.

\bibitem{SuiteSparse}
T.~A. Davis and Y.~Hu.
\newblock The University of Florida sparse matrix collection.
\newblock \emph{ACM Transactions on Mathematical Software}, 38(1):1--25, 2011.

\end{thebibliography}

\end{document}
